\section*{Chapitre \ref{ch:hf:short-horizon-alpha} --- Alpha à court horizon}

\begin{solution}{exo:hf:short-horizon-alpha:1}
Déséquilibre $(12-4)/16=0.5$~; prévision $0.30\times0.5=0.15$ ticks.
\end{solution}

\begin{solution}{exo:hf:short-horizon-alpha:2}
Un demi-tick plus les frais, 0.1 cent $=0.1$ tick~: 0.6 ticks.
\end{solution}

\begin{solution}{exo:hf:short-horizon-alpha:3}
Cours acheteur en hausse~: $+3$ (la nouvelle taille à l'achat, rien n'est retiré)~; cours vendeur inchangé en prix~: $-6+5=-1$. L'incrément vaut $3-1=2$ lots.
\end{solution}

\begin{solution}{exo:hf:short-horizon-alpha:4}
Le mid retarde sur le prix efficient de quelques secondes, de sorte qu'un mouvement récent tend à se poursuivre pendant que le carnet rattrape son retard. Sur une action réelle, où les cotations suivent
l'information plus vite, l'impact se résorbant en partie, le mouvement propre entre en général avec un coefficient faible ou négatif à ces horizons.
\end{solution}

\begin{solution}{exo:hf:short-horizon-alpha:5}
Déséquilibre (0.187) à une demi-seconde~; déséquilibre du flux d'ordres (0.372) à dix. Un teneur de marché qui maintient ses cotations pendant cinq secondes doit pondérer les variables les plus lentes~:
ajuster la prévision à l'horizon qu'il tient, et non au plus court.
\end{solution}

\begin{solution}{exo:hf:short-horizon-alpha:6}
La taille typique de la prévision est sa corrélation multipliée par la volatilité de la cible, bien plus petite que la cible elle-même~; seule une infime part des prévisions
dépasse le coût d'agression (0.01~\% dépassait 0.9 ticks ici). Sa valeur tient à de nombreuses petites décisions sur la cotation à afficher.
\end{solution}

\begin{solution}{exo:hf:short-horizon-alpha:7}
0.444, avec dix messages dont la médiane vaut 451 millisecondes.
\end{solution}

\begin{solution}{exo:hf:short-horizon-alpha:8}
Il s'exécute au mid, à l'instant du signal~: agresser le marché coûte la moitié de la fourchette et les frais, et le prix disponible un message plus tard a déjà
bougé avec le signal. Rejouer au cours vendeur (ou acheteur) du message suivant, frais inclus~; l'edge par transaction baisse d'au moins la demi-fourchette.
\end{solution}

\begin{solution}{pb:hf:short-horizon-alpha:1}
\begin{enumerate}
\item Une prévision sur des millisecondes à des secondes, à partir du carnet, des transactions et des instruments liés, servant surtout à décider quelles cotations afficher~; l'alpha intrajournalier
opère à l'échelle de la minute et paie le coût de l'agression.
\item Déséquilibre des files~; déséquilibre du flux d'ordres sur une seconde~; volume agressif net sur une seconde~; variation du mid du leader sur une seconde~; la
variation du mid propre sur une seconde.
\item 0.320, 0.285, 0.077, 0.288, 0.150~; ensemble 0.480.
\item Un carnet entièrement déséquilibré côté acheteur prévoit $+0.30$ ticks~; un mouvement du leader d'un tick, $+0.34$ ticks.
\item Le mid simulé reste en retard sur son prix efficient de quelques secondes (autocorrélation d'ordre un des rendements sur deux secondes, 0.44), ce qui rend le mouvement suivant inhabituellement
prévisible.
\item 0.28 à une demi-seconde, 0.37 à une seconde, 0.48 à deux, 0.54 à cinq, 0.50 à dix.
\item Environ deux variations moyennes de prix.
\item 0.477 (33 ms), 0.405 (0.9 s), 0.303 (2.3 s).
\item Les variations réalisées moyennes croissent avec la prévision d'un décile à l'autre, et suivent de près celle-ci aux extrêmes (décile supérieur~: prévision 0.31, réalisé 0.33).
\item Déplacer ou retirer le côté que la prévision annonce comme devant être servi~; en un tick, n'afficher que le côté favorable.
\item Agresser lorsque la prévision dépasse la demi-fourchette, plus les frais du preneur de liquidité, plus une marge~: $0.5+0.1+0.3=0.9$ ticks.
\item 0.01~\% des prévisions au-delà de 0.9 ticks, 5.8~\% au-delà de 0.3.
\item Aucun ajustement 16\,550 à 0.294~; skew 11\,050 à 0.451~; skew et \omterm{def:hf:short-horizon-alpha:take}{seuil d'agression} 10\,867 à 0.452~; un message de retard 10\,883 à 0.414.
\item 0.451 contre 0.294 ticks par action~; avec un message de retard, le mark-out passif baisse de 8~\% (de 0.452 à 0.414), tandis que le
coefficient d'information de la prévision baisse de moins de 1~\% (de 0.480 à 0.477).
\item Environ 25 agressions en six séances~: leur mark-out a une erreur type aussi grande que les écarts.
\item Les deux retirent un côté~; le fading utilise une prévision du mark-out de l'exécution, le skew une prévision de la direction du prix~; les deux abandonnent un tiers du
volume pour de meilleures exécutions.
\item Retour après balayage~: sur ce marché, le mouvement propre prédit la poursuite.
\item La durée de vie médiane des arbitrages est passée de 97 ms à 7 ms, tandis que le profit médian par opportunité est resté proche de 0.08 point d'indice.
\item Exécuter un message plus tard au prix alors disponible, frais inclus, et compter assez d'agressions (davantage de séances) pour une erreur type.
\item Décider lesquelles de ses propres cotations laisser au carnet d'ordres.
\end{enumerate}
\end{solution}

\begin{solution}{iq:hf:short-horizon-alpha:1}
Retirer ou relever le cours vendeur du fonds~: le mouvement du contrat à terme atteindra le fonds, et un cours vendeur périmé sera pris par quiconque lit le contrat à terme plus vite~; garder le
cours acheteur.
\iqlookfor{le mouvement du leader comme prévision et la protection du côté périmé.}
\end{solution}

\begin{solution}{iq:hf:short-horizon-alpha:2}
La variation de la file au cours acheteur (en comptant un meilleur cours acheteur comme sa nouvelle taille, un moins bon comme la perte de l'ancienne), moins la même grandeur pour la file au cours vendeur, sommée sur
les événements. Elle inclut les annulations et les nouveaux ordres à cours limité, qui déplacent le prix autant que les transactions, alors que le volume signé ne voit que les exécutions.
\iqlookfor{la définition et l'importance des annulations.}
\end{solution}

\begin{solution}{iq:hf:short-horizon-alpha:3}
Des cibles qui se chevauchent rendent les erreurs voisines corrélées~: utiliser des erreurs types HAC (livre~7, chapitre~6), évaluer sur des séances non utilisées pour
l'ajustement, et rapporter le coefficient d'information par horizon.
\iqlookfor{chevauchement, erreurs HAC, séances hors échantillon.}
\end{solution}

\begin{solution}{iq:hf:short-horizon-alpha:4}
Un seul moteur de variables pour les deux, alimenté message par message~; un test de rejeu qui fait tourner le moteur de production sur les flux stockés et compare sa sortie à
celle de la recherche, valeur par valeur.
\iqlookfor{une seule implémentation et un test de parité.}
\end{solution}

\begin{solution}{iq:hf:short-horizon-alpha:5}
Lorsque la prévision dépasse la demi-fourchette, plus les frais du preneur de liquidité, plus une marge~; la marge couvre l'erreur de la prévision, son âge au moment où l'ordre
arrive, et la probabilité que des traders plus rapides aient pris le mouvement en premier.
\iqlookfor{le seuil et ses composantes.}
\end{solution}

\begin{solution}{iq:hf:short-horizon-alpha:6}
La moyenne conditionnelle a pour écart type $\rho s$, donc $P(|\hat\alpha|>c)=2\Phi(-c/(\rho s))$~: avec $\rho s=0.288$ et $c=0.9$, $2\Phi(-3.13)
\approx0.18~\%$.
\iqlookfor{la variance d'une prévision et la queue gaussienne.}
\end{solution}
